% Part: first-order-logic
% Chapter: syntax-and-semantics
% Section: free-vars-sentences

\documentclass[../../../include/open-logic-section]{subfiles}

\begin{document}

\olfileid{fol}{syn}{fvs}

\olsection{Vrije \printtoken{P}{variable} en \printtoken{P}{sentence}}

\begin{defn}[Vrije voorkomens van !!a{variable}]
\ollabel{defn:free-occ}
De \emph{vrije} voorkomens van !!a{variable} in !!a{formula} worden
als volgt inductief gedefinieerd:
\begin{enumerate}
\item \indcase*{!A}{$!A$ atomair is}{alle voorkomens van
  !!{variable}s in $\indfrm$ vrij zijn.}

\tagitem{prvNot}{\indcase{!A}{\lnot !B}{de vrije voorkomens van
  !!{variable}s in $\indfrm$ precies die in $!B$ zijn.}}{}

\item \indcase{!A}{(!B \ast !C)}{de vrije voorkomens van
  !!{variable}s in $\indfrm$ bestaan uit de vrije voorkomens in $!B$
  samen met die in~$!C$.}

\tagitem{prvAll}{\indcase{!A}{\lforall[x][!B]}{de vrije voorkomens van
  !!{variable}s in $\indfrm$ alle vrije voorkomens in~$!B$ zijn,
  behalve de voorkomens van~$x$.}}{}

\tagitem{prvEx}{\indcase{!A}{\lexists[x][!B]}{de vrije voorkomens van
  !!{variable}s in $\indfrm$ alle vrije voorkomens in~$!B$ zijn,
  behalve de voorkomens van~$x$.}}{}
\end{enumerate}
\end{defn}

\begin{defn}[Gebonden variabelen]
Een voorkomen van !!a{variable} in een formule~$!A$ is \emph{gebonden}
als het niet vrij is.
\end{defn}

\begin{prob}
Geef naar het voorbeeld van \olref[fol][syn][fvs]{defn:free-occ} een
inductieve definitie van de gebonden voorkomens van variabelen.
\end{prob}

\begin{defn}[Bereik]
\iftag{prvAll}{Als $\lforall[x][!B]$ als deelformule voorkomt in een
  formule~$!A$, dan heet het overeenkomstige voorkomen van~$!B$ in~$!A$
  het \emph{bereik} van het overeenkomstige voorkomen van~$\lforall[x]$.
  \iftag{prvEx}{Hetzelfde geldt voor $\lexists[x]$.}{}}{Als
  $\lexists[x][!B]$ als deelformule voorkomt in een formule~$!A$, dan
  heet het overeenkomstige voorkomen van~$!B$ in~$!A$ het \emph{bereik}
  van het overeenkomstige voorkomen van~$\lexists[x]$.}

Als $!B$ het bereik is van een voorkomen van de kwantor
\iftag{prvAll}{$\lforall[x]$\iftag{prvEx}{ of
    $\lexists[x]$}{}}{$\lexists[x]$} in~$!A$, dan zijn de vrije
voorkomens van $x$ in~$!B$ gebonden in
\iftag{prvAll}{$\lforall[x][!B]$\iftag{prvEx}{ en
$\lexists[x][!B]$}{}}{$\lexists[x][!B]$}. We zeggen dat deze voorkomens
\emph{worden gebonden door} het genoemde voorkomen van de kwantor.
\end{defn}

\begin{ex}
Beschouw de volgende formule:
\[
\lexists[\Obj v_0][\underbrace{\Atom{\Obj A^2_0}{\Obj v_0,\Obj v_1}}_{!B}] 
\]
$!B$ is het bereik van $\lexists[\Obj v_0]$.
De kwantor bindt het voorkomen van $\Obj v_0$ in $!B$, maar niet het
voorkomen van $\Obj v_1$. $\Obj v_1$ is hier dus een vrije variabele.


We kunnen nu zien hoe dit werkt in een ingewikkelder
!!{formula}~$!A$:
\[
\lforall[\Obj v_0][\underbrace{(\Atom{\Obj A^1_0}{\Obj v_0} \lif
    \Atom{\Obj A^2_0}{\Obj v_0, \Obj v_1})}_{!B}] \lif \lexists[\Obj
  v_1][\underbrace{(\Atom{\Obj A^2_1}{\Obj v_0, \Obj v_1} \lor \lforall[\Obj v_0][\overbrace{\lnot \Atom{\Obj A^1_1}{\Obj v_0}}^{!D}])}_{!C}]
\]
$!B$ is het bereik van de eerste $\lforall[\Obj v_0]$, $!C$ is het bereik
van $\lexists[\Obj v_1]$ en $!D$ is het bereik van de tweede
$\lforall[\Obj v_0]$. De eerste $\lforall[\Obj v_0]$ bindt de voorkomens
van $\Obj v_0$ in~$!B$, $\lexists[\Obj v_1]$ bindt het voorkomen van
$\Obj v_1$ in $!C$ en de tweede $\lforall[\Obj v_0]$ bindt het voorkomen
van $\Obj v_0$ in~$!D$. Het eerste voorkomen van $\Obj v_1$ en het vierde
voorkomen van $\Obj v_0$ zijn vrij in~$!A$. Het laatste voorkomen van
$\Obj v_0$ is vrij in $!D$, maar gebonden in $!C$ en~$!A$.
\end{ex}

\begin{defn}[Zin]
!!^a{formula}~$!A$ is \article{sentence} \emph{!!{sentence}} dan en
slechts dan als de formule geen vrije voorkomens van !!{variable}s bevat.
\end{defn}

% add examples!
\end{document}
